\documentclass[11pt,a4paper]{article}
\usepackage[lmargin=1.0in,rmargin=1.0in,bottom=1.0in,top=1.0in,twoside=False]{geometry}


\usepackage{amsmath,amsthm}
\usepackage{mathtools}
\usepackage{graphicx}%,algorithmic,algorithm}%, verbatim}
\usepackage{subcaption}
\usepackage{enumerate}
\usepackage{tikz}
\usetikzlibrary{shapes}
\usetikzlibrary{arrows}
\usetikzlibrary{calc}

\tikzstyle{node}=[fill=black, shape=circle, inner sep=2pt]
\tikzstyle{node-is}=[fill={rgb,255: red,191; green,0; blue,64}, shape=circle, inner sep=2pt]
\tikzstyle{none}=[inner sep=0pt]
\tikzstyle{directed}=[->]
\tikzstyle{undirected}=[-]
\tikzstyle{bi-directional}=[<->]

\usepackage[T1]{fontenc}

\usepackage{enumitem}

\usepackage{microtype}
\usepackage{hyperref}
\usepackage{amsfonts}
\usepackage{comment}
\usepackage{mathrsfs}
\usepackage[capitalize, nameinlink]{cleveref}
\usepackage{todonotes}
%\usepackage[ocgcolorlinks, linkcolor={blue}, citecolor={brown}]{hyperref}
%\usepackage{cleveref}

\usepackage{stmaryrd}

\newtheorem{lemma}{Lemma}[section]
\newtheorem{theorem}[lemma]{Theorem}
\newtheorem{corollary}[lemma]{Corollary}
\newtheorem{definition}[lemma]{Definition}
%\theoremstyle{definition}
%\newtheorem{definition}[lemma]{Definition}
\newtheorem{claim}[lemma]{Claim}
\newtheorem{prop}[lemma]{Proposition}
\newtheorem{obs}[lemma]{Observation}

\renewcommand{\mod}{\ \mathsf{mod}\ }
\newcommand{\Oh}{\mathcal{O}}
\newcommand{\sym}{\mathcal{S}}
\newcommand{\cA}{\mathcal{A}}
\newcommand{\cB}{\mathcal{B}}

\newcommand{\kacper}[2][]{\todo[color=green!20,inline,#1]{{\textbf{Kacper:} #2}}}


\begin{document}
\iftrue
\author{
  Kacper Kluk% 
  \thanks{Institute of Informatics, University of Warsaw, Poland. \texttt{k.kluk@uw.edu.pl}.}
  \and%
  Jesper Nederlof%
  \thanks{Department of Information and Computing Sciences, Utrecht University, The Netherlands. \texttt{j.nederlof@uu.nl}.}
}
\else
\author{}
\fi

\title{Lower bounds on pure dynamic programming for connectivity problems on graphs of bounded path-width
\iftrue
  \thanks{
  K.K. is supported by Polish National Science Centre SONATA BIS-12 grant number 2022/46/E/ST6/00143. J.N is supported by the project COALESCE that has received funding from the European Research Council (ERC), grant agreement No 853234.}\fi}

\date{}
\maketitle

\begin{abstract}
We give unconditional parameterized complexity lower bounds on pure dynamic programming algorithms -- as modeled by \emph{tropical circuits} -- for connectivity problems such as the Traveling Salesperson Problem. Our lower bounds are higher than the currently fastest algorithms that rely on algebra and give evidence that these algebraic aspects are unavoidable for competitive worst case running times. 

Specifically, we study input graphs with a small width parameter such as treewidth and pathwidth and show that for any $k$ there exists a graph $G$ of pathwidth at most $k$ and $k^{O(1)}$ vertices such that any tropical circuit calculating the optimal value of a Traveling Salesperson round tour uses at least $2^{\Omega(k \log \log k)}$ gates. We establish this result by linking tropical circuit complexity to the nondeterministic communication complexity of specific compatibility matrices. These matrices encode whether two partial solutions combine into a full solution, and Raz and Spieker~[Combinatorica 1995] previously proved a lower bound for this complexity measure.
\end{abstract}



\input{intro}
\input{prelims}
\input{mwis}
\input{comcom}
\input{tsp}
\input{st}


\section{Conclusion}
\label{sec:con}

As noted earlier, the exact nondeterministic communication complexity of the matchings compatibility matrix is still open: While $C_k =k^{O(k)}$ holds trivially, the currently best known lower bound is $C_k = 2^{\Omega(k \log \log k)}$~\cite{raz1995log}. 
Proving this lower bound would automatically imply the same bound for \Cref{thm:main_dtsp}, \Cref{thm:main_tsp}, and \Cref{thm:main_dst}. In particular, this complexity would be asymptotically tight (up to factors polynomial in $N$), as tropical circuits of size $2^{\Oh(k \log k)} \cdot N^{\Oh(1)}$ can be easily constructed for all these problems by following algorithms based on naive dynamic programming on path decompositions.
In the reverse direction, better upper bounds on $C_k$ seem also useful for obtaining better tropical circuits for $DTSP_G,TSP_G$ and $DST_G$, and possibly even for obtaining faster algorithm for TSP parameterized by pathwidth. It seems however, as also expressed by the authors of~\cite{raz1995log}, that $C_k$ is closer to $2^{\Omega(k \log k)}$ than it is to $2^{\Theta(k \log \log k)}$.


Another natural opportunity for further work would be to find lower bounds for other computational problems, for example it seems plausible that with tools from Theorem~\ref{thm:main_is} and the reduction ideas from~\cite{LokshtanovMS18b} one can also obtain a $\Omega(3^k)$ lower bound for tropical circuits calculating the minimum weight dominating set of a graph with pathwidth $k$.

More ambitiously, a natural open question is whether the lower bounds on tropical circuits as defined in this paper can be generalized to more expressive variants of tropical circuits (see the book by Jukna~\cite{Jukna2023}).

\bibliographystyle{alpha}
\bibliography{ref}

\end{document}